\documentclass[10pt,a4paper]{amsart}
\usepackage[margin=19mm]{geometry}
\usepackage{amsmath,amssymb,mathtools}
\usepackage[scr=rsfs,cal=euler]{mathalfa}
\usepackage{xfrac}
\usepackage[unicode,hidelinks,bookmarks=false]{hyperref}
\newcommand{\ie}{i.e.\ }
\newcommand{\cf}{cf.\ }
\newcommand{\resp}{respectively\ }
\newcommand{\Subsection}{\subsection}
\providecommand{\nobreakdash}{\nobreak\hbox{-}}
\setlength{\emergencystretch}{2em}
\multlinegap0pt
\usepackage{bm}
\usepackage{bbm}
\usepackage{dsfont}
\usepackage{xspace}
\usepackage{cancel}
\usepackage{mathtools}
\usepackage{amsthm}
\makeatletter
\@ifundefined{theorem}{\newtheorem{theorem}{Theorem}}{}
\makeatother
\newcommand{\Ft}{\widetilde{F}}
\newcommand{\Mbarzn}{\overline{\mathcal M}_{0,n}}
\newcommand{\Mbar}{\overline{\mathcal M}}
\newcommand{\PP}{\mathbb{P}} 
\newcommand{\cM}{\mathcal{M}} 
\title[Stratified motivic invariants and bivariate deformations of ]{Stratified motivic invariants and bivariate deformations of Poincar\'e polynomials}
\author{Gergely B\'erczi, Young-Hoon Kiem}
\date{Source: arXiv:2607.23890v1 (CC BY 4.0)}
\begin{document}
\maketitle

\noindent\emph{Source and licence.} Stratified motivic invariants and bivariate deformations of Poincar\'e polynomials. Version arXiv:2607.23890v1 (CC BY 4.0), distributed under the Creative Commons Attribution 4.0 International licence, as recorded in the arXiv licence metadata for this exact version and shown on the abstract page https://arxiv.org/abs/2607.23890v1. Full rights notice: licenses/arxiv-2607.23890-CC-BY-4.0.txt.\par\smallskip

\noindent\textit{Editorial scope.} Counted unit: the Theorem whose source label is \texttt{prop:FM-coefficient-interpretation} in the article (source0.tex lines 726--744, source label \texttt{prop:FM-coefficient-interpretation}), delivered together with the article's own complete original proof of it (source0.tex lines 746--811, one \texttt{proof} environment of 66 lines). No mathematics is written, repaired or re-derived: statement and proof are the article's own text line for line. Required context retained uncounted: 10 (lines 1108--1111). Every definition, display and label of those lines is retained unchanged. Omitted from this package and not redistributed: 1--725, 745--745, 812--1107, 1112--1123. Nothing inside a retained excerpt is omitted or altered: every retained body is delivered exactly as the article's lines are written, so no omission receipt of any kind accompanies this package. Citations: the retained excerpts carry no citation command at all, so no bibliography entry of the article is transcribed and no reference is redistributed. Reference adaptation: none was needed and none was made; every reference inside the delivered unit resolves inside it, so no unresolved reference or citation is produced. Printed slips kept literal: no printed word, symbol, hypothesis or number of the counted statement or of its proof was altered. Layout and class: the article's own class is \texttt{amsart} with packages including geometry, amsmath, amssymb, amsthm, mathtools, enumitem, microtype, hyperref, tikz, pgfplots; the wrapper is the lane's pinned \texttt{amsart} editorial wrapper at 10pt on A4, which restates the theorem-like environments the retained text needs and adds the article's own macro definitions that the retained text needs (\texttt{\textbackslash Ft}, \texttt{\textbackslash Mbar}, \texttt{\textbackslash Mbarzn}, \texttt{\textbackslash PP}, \texttt{\textbackslash cM}), with extra wrapper packages (bm, bbm, dsfont, xspace, cancel, mathtools). The delivered numbering is the unit's own. Assets: no retained span contains a figure environment, a graphics-inclusion command, a PDF-page inclusion, a table or a TikZ picture; no image file is copied into this package, the package declares no asset dependency and no image file is redistributed with it. Licence: version arXiv:2607.23890v1 of this article is distributed under the Creative Commons Attribution 4.0 International licence, as recorded in the arXiv licence metadata for this exact version and shown on the abstract page https://arxiv.org/abs/2607.23890v1. A full rights notice is delivered at licenses/arxiv-2607.23890-CC-BY-4.0.txt. Characters: every retained excerpt is pure ASCII, so the delivered engine, XeLaTeX with its default Latin Modern text font, reports no missing character.
\begin{theorem}\label{10}
$$F_m(y,t)=\sum_{v\ge 1} y^v\,\PP_{\sqrt{t}}(\cM_{0,m+1,v}), \quad [y^v]F_m(y,u^2)=\PP_{u}(\cM_{0,m+1,v})$$ where $\PP_u(\cM_{0,m+1,v})$ denotes the virtual Poincar\'e polynomial of the union $\cM_{0,m+1,v}$
of strata of stable curves in $\Mbarzn$ with precisely $v$ irreducible components.  
\end{theorem}

\begin{theorem}[Coefficient interpretation for the FM deformation]
\label{prop:FM-coefficient-interpretation}
For every $n\ge 1$ and every $0\le q\le n-1$, we have
\begin{equation}\label{eq:FM-coeff-full}
  [y^q]\Ft_n(y,u^2)=\PP_u(X_{n,q})
\end{equation}
and
\begin{equation}\label{eq:FM-coeff-residual}
  [y^q]\widehat F_n(y,u^2)=\PP_u(X_{n,q}^{\infty}).
\end{equation}
In particular,
\[
  \widehat F_n(1,u^2)=\PP_u\bigl(\mathrm{ev}_1^{-1}(\infty)\bigr).
\]
Thus the residual deformation has the following modular meaning: the
coefficient of $y^q$ is the virtual Poincar\'e polynomial of the locus of
stable maps whose source curve has precisely $q$ contracted components, after
fixing the image of the first marking to be $\infty$.
\end{theorem}

\begin{proof}
Let $\Pi_n$ be the set of set partitions of $\{1,\ldots,n\}$.  Given
$f\in X_n$, two markings $p_i,p_j$ determine the same block precisely when
their images on the degree-one component coincide.  Thus a partition
\[
  \pi=\{B_1,\ldots,B_k\}\in \Pi_n
\]
records the $k$ support points on the component $C_0\cong \PP^1$.
For a block $B$ of size $m$, the contracted tree attached over the
corresponding support point is a rooted stable rational curve with $m$ labelled
markings and one additional root marking, hence is parametrized by
$\Mbar_{0,m+1}$.  When $m=1$ there is no contracted component over that
support point, and this agrees with the convention $F_1(y,t)=1$.
For $m\ge 2$, the interpretation \eqref{10} says that the coefficient of
$y^v$ in $F_m(y,u^2)$ is the virtual Poincar\'e polynomial of the locus of
rooted bubble trees with exactly $v$ irreducible components.  These components
are precisely the contracted components of the stable map.

For a fixed partition $\pi$ with $k$ blocks, the support points form the
ordered configuration space of $k$ distinct points of $\PP^1$, ordered by the
blocks of $\pi$.  Its virtual Poincar\'e polynomial is
\[
  \PP_u(\operatorname{Conf}_k(\PP^1))
  =(t+1)t(t-1)\cdots(t-k+2),
  \qquad t=u^2.
\]
Therefore the contribution of all stable maps with support partition $\pi$ is
\[
  (t+1)t(t-1)\cdots(t-k+2)\prod_{B\in \pi}F_{|B|}(y,t).
\]
Summing over all partitions gives
\begin{equation}\label{eq:FM-partition-full}
  \sum_{q\ge 0}y^q\PP_u(X_{n,q})
  =
  \sum_{\pi\in\Pi_n}
  (t+1)t(t-1)\cdots(t-|\pi|+2)
  \prod_{B\in\pi}F_{|B|}(y,t).
\end{equation}
By the exponential formula, the right hand side is exactly the coefficient of
$x^n/n!$ in $(1+\Phi(x,y,t))^{t+1}$.  This proves
\eqref{eq:FM-coeff-full}.

For the residual statement, fix $\mathrm{ev}_1=\infty$.  If $B_0$ is the
block of $\pi$ containing the label $1$, then the support point corresponding
to $B_0$ is forced to be $\infty$.  The remaining $k-1$ support points form
the ordered configuration space of $k-1$ distinct points of
$\mathbb A^1=\PP^1\setminus\{\infty\}$.  Hence
\[
  \PP_u(\operatorname{Conf}_{k-1}(\mathbb A^1))
  =t(t-1)\cdots(t-k+2),
\]
with the convention that this product is $1$ when $k=1$.  Thus
\begin{equation}\label{eq:FM-partition-residual}
  \sum_{q\ge 0}y^q\PP_u(X_{n,q}^{\infty})
  =
  \sum_{\pi\in\Pi_n}
  t(t-1)\cdots(t-|\pi|+2)
  \prod_{B\in\pi}F_{|B|}(y,t).
\end{equation}
Again by the exponential formula, the right hand side is the coefficient of
$x^n/n!$ in
\[
  \frac{(1+\Phi(x,y,t))^{t+1}-1}{t+1}.
\]
This proves \eqref{eq:FM-coeff-residual}.
\end{proof}

\end{document}
